\documentclass[11pt]{ctexart}
\usepackage[a4paper,margin=1in]{geometry}
\usepackage{longtable,booktabs}
\title{Geant4-Agent 回归测试报告（中文）}
\author{自动化评测流程}
\date{2026-03-07\_142103}
\begin{document}
\maketitle
\section{统一 Pass/Fail 总表}
\begin{longtable}{p{0.06\linewidth}p{0.08\linewidth}p{0.07\linewidth}p{0.08\linewidth}p{0.07\linewidth}p{0.07\linewidth}p{0.07\linewidth}p{0.07\linewidth}p{0.09\linewidth}p{0.06\linewidth}p{0.21\linewidth}}
\toprule
ID & 领域 & 风格 & 参数完整性 & 期望初始 & 实际初始 & 期望最终 & 实际最终 & 缺参(初->终) & 结果 & 失败原因 \\
\midrule
MTM01 & medical & fuzzy & missing & 失败 & 失败 & 通过 & 通过 & 13->0 & 通过 & - \\
MTM02 & aerospace & precise & missing & 失败 & 失败 & 通过 & 通过 & 2->0 & 通过 & - \\
MTM03 & industrial\_ndt & precise & missing & 失败 & 失败 & 通过 & 通过 & 12->0 & 通过 & - \\
MTM04 & physics & precise & missing & 失败 & 失败 & 通过 & 通过 & 8->0 & 通过 & - \\
MTM05 & nuclear\_engineering & fuzzy & missing & 失败 & 失败 & 通过 & 通过 & 5->0 & 通过 & - \\
MTM06 & security\_screening & fuzzy & missing & 失败 & 失败 & 通过 & 通过 & 12->0 & 通过 & - \\
MTM07 & semiconductor & precise & missing & 失败 & 失败 & 通过 & 通过 & 3->0 & 通过 & - \\
MTM08 & education & fuzzy & missing & 失败 & 失败 & 通过 & 通过 & 7->0 & 通过 & - \\
MTM09 & environmental\_radiation & fuzzy & missing & 失败 & 失败 & 通过 & 通过 & 3->0 & 通过 & - \\
MTM10 & high\_energy\_physics & precise & missing & 失败 & 失败 & 通过 & 通过 & 10->0 & 通过 & - \\
MTM11 & materials\_update & precise & full & 通过 & 通过 & 通过 & 通过 & 0->0 & 通过 & - \\
MTM12 & overwrite\_reject & precise & full & 通过 & 通过 & 通过 & 通过 & 0->0 & 通过 & - \\
\bottomrule
\end{longtable}
\section{基准对比}
\begin{itemize}
\item pass\_rate=1.000
\item initial\_alignment\_rate=1.000
\item final\_alignment\_rate=1.000
\item avg\_turns\_per\_case=2.583
\item missing\_reduction\_rate=1.000
\item confirm\_apply\_success\_rate=1.000
\item rollback\_after\_confirm\_rate=0.000
\item internal\_leak\_turn\_rate=0.065
\item repeated\_assistant\_case\_rate=0.000
\item stale\_confirm\_case\_rate=0.000
\item delta\_pass\_rate=+0.050
\item regressions=无
\item improvements=无
\end{itemize}
\appendix
\section{附录：完整对话与内部处理链（自然语言）}
\subsection{案例 MTM01（medical）}
\textbf{对话：}
\begin{itemize}
\item 用户：我先要一个用于医学成像的初步辐照场景，细节还没定，你先告诉我还缺什么。
\item 系统：你提到的医学成像场景，具体需要什么样的几何形状呢？比如是简单的盒体、圆柱，还是更复杂的结构？另外，如果选择盒体，它在x方向上的尺寸大概是多少？
\item 用户：采用 300 mm x 300 mm x 300 mm 的水箱，材料 G4\_WATER；源用束流 gamma，能量 6 MeV，位置 (0,0,-400) mm，方向 +z；物理列表 QBBC；输出 root。
\item 系统：配置已完成。当前已确认的设置有：物理列表的选择依据和来源，以及材料的选择依据和来源。
\end{itemize}
\textbf{内部链路（自然语言）：}
\begin{itemize}
\item 第 1 轮：动作=ask\_clarification，阶段=geometry，闭环状态=False。
\item 缺失字段：geometry.structure、geometry.params.module\_x、geometry.params.module\_y、geometry.params.module\_z、materials.volume\_material\_map、source.type、source.particle、source.energy、source.position、source.direction、physics.physics\_list、output.format、output.path。
\item LLM 使用=True；回退=无；后端=llm\_semantic\_frame+runtime\_semantic。
\item 第 2 轮：动作=finalize，阶段=materials，闭环状态=True。
\item 缺失字段：无。
\item LLM 使用=True；回退=无；后端=llm\_slot\_frame+runtime\_semantic。
\end{itemize}
\subsection{案例 MTM02（aerospace）}
\textbf{对话：}
\begin{itemize}
\item 用户：I need a first-pass proton irradiation setup for a silicon device; keep it simple for now.
\item 系统：Geometry: updated geometry structure, box size along x, box size along y. Materials: updated materials.selected\_materials, volume-to-material binding, material selection source. Source: updated source type, particle type, source energy. Physics: updated physics list, physics.backup\_physics\_list, physics selection reasons. Output: still needs output format, output path.
\item 用户：Use a silicon box 50 mm x 50 mm x 5 mm, beam source proton 100 MeV from (0,0,-100) mm along +z, physics FTFP\_BERT, output json.
\item 系统：This would overwrite an already confirmed value.

Pending changes:
- Box size along x: 1000.0 → 50.0
- Box size along y: 1000.0 → 50.0

Reply **'confirm'** to apply these changes, or say **'keep existing'** to keep the current values.
\item 用户：confirm
\item 系统：Your configuration is complete. The following items have been confirmed: geometry program, source energy, source type, and the box size along the Z-axis.
\end{itemize}
\textbf{内部链路（自然语言）：}
\begin{itemize}
\item 第 1 轮：动作=summarize\_progress，阶段=materials，闭环状态=False。
\item 缺失字段：output.format、output.path。
\item LLM 使用=True；回退=无；后端=llm\_slot\_frame+runtime\_semantic。
\item 第 2 轮：动作=confirm\_overwrite，阶段=source，闭环状态=False。
\item 缺失字段：无。
\item LLM 使用=True；回退=无；后端=llm\_slot\_frame+runtime\_semantic。
\item 第 3 轮：动作=finalize，阶段=physics，闭环状态=True。
\item 缺失字段：无。
\item LLM 使用=True；回退=无；后端=llm\_semantic\_frame+runtime\_semantic+confirmed\_pending\_overwrite。
\end{itemize}
\subsection{案例 MTM03（industrial\_ndt）}
\textbf{对话：}
\begin{itemize}
\item 用户：我要做一个 gamma 穿透铁靶的小测试，但我还没把完整参数说完。
\item 系统：目前我们已经确认了材料和源的基本设置，包括使用铁作为靶材，以及选择 gamma 作为粒子类型。接下来还需要补充几何结构的具体尺寸、体积与材料的绑定关系，以及源的详细参数和物理列表等信息。请继续提供这些细节，以便我帮你完成配置。
\item 用户：请设为半径 40 mm、半长 80 mm 的铁圆柱，材料 G4\_Fe；点源 gamma，能量 2 MeV，位置 (0,0,-200) mm，方向 +z；物理列表 Shielding；输出 csv。
\item 系统：配置已完成。当前已确认的参数包括：物理列表选择依据、物理列表选择来源、体积与材料绑定、几何根体积名。
\end{itemize}
\textbf{内部链路（自然语言）：}
\begin{itemize}
\item 第 1 轮：动作=summarize\_progress，阶段=geometry，闭环状态=False。
\item 缺失字段：geometry.structure、geometry.params.module\_x、geometry.params.module\_y、geometry.params.module\_z、materials.volume\_material\_map、source.type、source.energy、source.position、source.direction、physics.physics\_list、output.format、output.path。
\item LLM 使用=True；回退=无；后端=llm\_slot\_frame+runtime\_semantic。
\item 第 2 轮：动作=finalize，阶段=materials，闭环状态=True。
\item 缺失字段：无。
\item LLM 使用=True；回退=无；后端=llm\_slot\_frame+runtime\_semantic。
\end{itemize}
\subsection{案例 MTM04（physics）}
\textbf{对话：}
\begin{itemize}
\item 用户：Gamma point source 1 MeV at (0,0,-100) toward +z; please build the rest.
\item 系统：Geometry: still needs geometry structure, box size along x. Materials: still needs volume-to-material binding. Source: updated source type, particle type, source energy. Physics: still needs physics list. Output: still needs output format, output path.
\item 用户：Set copper box 1 m x 1 m x 1 m, physics FTFP\_BERT, output hdf5.
\item 系统：Your configuration is complete. The physics and material selections have been confirmed.
\end{itemize}
\textbf{内部链路（自然语言）：}
\begin{itemize}
\item 第 1 轮：动作=summarize\_progress，阶段=geometry，闭环状态=False。
\item 缺失字段：geometry.structure、geometry.params.module\_x、geometry.params.module\_y、geometry.params.module\_z、materials.volume\_material\_map、physics.physics\_list、output.format、output.path。
\item LLM 使用=True；回退=无；后端=llm\_slot\_frame+runtime\_semantic。
\item 第 2 轮：动作=finalize，阶段=materials，闭环状态=True。
\item 缺失字段：无。
\item LLM 使用=True；回退=无；后端=llm\_slot\_frame+runtime\_semantic。
\end{itemize}
\subsection{案例 MTM05（nuclear\_engineering）}
\textbf{对话：}
\begin{itemize}
\item 用户：我想做一个中子屏蔽的概念验证，现在先给一个缺参开头。
\item 系统：好的，我们目前已经完成了不少配置：

*   **几何**：已确定使用盒体结构。
*   **材料**：已选定材料并绑定到几何体。
*   **源**：已确定使用中子源，并设置了能量。
*   **物理**：已选定物理列表。
*   **输出**：已设置好输出格式和路径。

接下来，还需要您补充一些尺寸和源的具体信息：
*   盒体的 **x、y、z 方向尺寸**。
*   中子源的 **位置** 和 **方向**。

请告诉我这些参数，我们就可以继续了。
\item 用户：使用 200 mm x 200 mm x 200 mm 的铅块，材料 G4\_Pb；点源 neutron，能量 14 MeV，位置 (0,0,-50) mm，方向 +z；物理列表 QGSP\_BERT\_HP；输出 xml。
\item 系统：我注意到新的设置会覆盖之前已经确认过的值，具体是：

1.  材料选择：从 **铜 (G4\_Cu)** 改为 **铅 (G4\_Pb)**。
2.  粒子类型：从 **伽马 (gamma)** 改为 **中子 (neutron)**。

如果您确认应用这些修改，请直接回复 **“确认”**；如果您想保留原来的值，请直接回复 **“保持原值”**。
\item 用户：确认
\item 系统：好的，配置已全部完成。目前我们确认的设置包括：体积与材料绑定、输出路径、输出格式和物理列表。
\end{itemize}
\textbf{内部链路（自然语言）：}
\begin{itemize}
\item 第 1 轮：动作=summarize\_progress，阶段=geometry，闭环状态=False。
\item 缺失字段：geometry.params.module\_x、geometry.params.module\_y、geometry.params.module\_z、source.position、source.direction。
\item LLM 使用=True；回退=无；后端=llm\_slot\_frame+runtime\_semantic。
\item 第 2 轮：动作=confirm\_overwrite，阶段=materials，闭环状态=False。
\item 缺失字段：无。
\item LLM 使用=True；回退=无；后端=llm\_slot\_frame+runtime\_semantic。
\item 第 3 轮：动作=finalize，阶段=source，闭环状态=True。
\item 缺失字段：无。
\item LLM 使用=True；回退=无；后端=llm\_semantic\_frame+runtime\_semantic+confirmed\_pending\_overwrite。
\end{itemize}
\subsection{案例 MTM06（security\_screening）}
\textbf{对话：}
\begin{itemize}
\item 用户：Need a quick cargo-screening toy setup with forward photons, but I have not fixed the target shape, material, source energy, or output format yet.
\item 系统：Here's a summary of your setup so far and what's still needed.

**Completed:**
- **Source:** Particle type (photon), source selection method, and reasons for selection are set.

**Still Required:**
- **Geometry:** Shape/structure and box dimensions (X, Y, Z).
- **Materials:** Assignment of material to your volume.
- **Source:** Source type and energy.
- **Physics:** Physics list selection.
- **Output:** Output format and file path.
\item 用户：Use steel cylinder target radius 60 mm half-length 120 mm, beam gamma 3 MeV from (0,0,-500) mm to +z, physics Shielding, output root.
\item 系统：Your configuration is complete. The physics and material selections have been confirmed.
\end{itemize}
\textbf{内部链路（自然语言）：}
\begin{itemize}
\item 第 1 轮：动作=summarize\_progress，阶段=geometry，闭环状态=False。
\item 缺失字段：geometry.structure、geometry.params.module\_x、geometry.params.module\_y、geometry.params.module\_z、materials.volume\_material\_map、source.type、source.energy、source.position、source.direction、physics.physics\_list、output.format、output.path。
\item LLM 使用=True；回退=无；后端=llm\_slot\_frame+runtime\_semantic。
\item 第 2 轮：动作=finalize，阶段=materials，闭环状态=True。
\item 缺失字段：无。
\item LLM 使用=True；回退=无；后端=llm\_slot\_frame+runtime\_semantic。
\end{itemize}
\subsection{案例 MTM07（semiconductor）}
\textbf{对话：}
\begin{itemize}
\item 用户：我想搭一个硅薄片的中子照射测试，但先只给一个简述。
\item 系统：目前我们已经完成了几何、材料和源的配置，接下来还需要你补充物理列表、输出格式和输出路径这三项信息。
\item 用户：使用 20 mm x 20 mm x 1 mm 的硅片，材料 G4\_Si；各向同性 neutron，能量 2 MeV，位置 (0,0,0) mm；物理列表 QGSP\_BERT\_HP；输出 json。
\item 系统：我注意到新的设置会覆盖之前已确认的参数。具体修改如下：

- 盒体 x 方向尺寸：从 1000.0 改为 20.0
- 盒体 y 方向尺寸：从 1000.0 改为 20.0

如果你确认应用这些修改，请回复“确认”；如果想保留原来的数值，请回复“保持原值”。
\item 用户：确认
\item 系统：配置已完成。当前已确认的参数包括：源位置、源能量、源类型以及盒体的 z 方向尺寸。
\end{itemize}
\textbf{内部链路（自然语言）：}
\begin{itemize}
\item 第 1 轮：动作=summarize\_progress，阶段=materials，闭环状态=False。
\item 缺失字段：physics.physics\_list、output.format、output.path。
\item LLM 使用=True；回退=无；后端=llm\_slot\_frame+runtime\_semantic。
\item 第 2 轮：动作=confirm\_overwrite，阶段=source，闭环状态=False。
\item 缺失字段：无。
\item LLM 使用=True；回退=无；后端=llm\_slot\_frame+runtime\_semantic。
\item 第 3 轮：动作=finalize，阶段=physics，闭环状态=True。
\item 缺失字段：无。
\item LLM 使用=True；回退=无；后端=llm\_semantic\_frame+runtime\_semantic+confirmed\_pending\_overwrite。
\end{itemize}
\subsection{案例 MTM08（education）}
\textbf{对话：}
\begin{itemize}
\item 用户：For classroom demo, I want gamma attenuation in copper; you can ask only essential missing parts.
\item 系统：Geometry: still needs geometry structure, box size along x. Materials: updated materials.selected\_materials, material selection source, material selection reasons; still needs volume-to-material binding. Source: updated particle type, source energy, source position. Physics: updated physics list, physics.backup\_physics\_list, physics selection reasons. Output: still needs output format, output path.
\item 用户：Okay, copper box 100 mm x 100 mm x 100 mm, point gamma 1 MeV at (0,0,-200) mm, direction +z, physics FTFP\_BERT, output csv.
\item 系统：This would overwrite an already confirmed value.

Pending changes:
- Source position: from [0.0, 0.0, -100.0] to [0.0, 0.0, -200.0]
- Physics list: from Shielding to FTFP\_BERT

Reply "confirm" to apply these changes, or say "keep existing" to keep the current values.
\item 用户：confirm
\item 系统：Configuration complete. The physics list, source position, output path, and output format are now confirmed.
\end{itemize}
\textbf{内部链路（自然语言）：}
\begin{itemize}
\item 第 1 轮：动作=summarize\_progress，阶段=geometry，闭环状态=False。
\item 缺失字段：geometry.structure、geometry.params.module\_x、geometry.params.module\_y、geometry.params.module\_z、materials.volume\_material\_map、output.format、output.path。
\item LLM 使用=True；回退=无；后端=llm\_slot\_frame+runtime\_semantic。
\item 第 2 轮：动作=confirm\_overwrite，阶段=materials，闭环状态=False。
\item 缺失字段：无。
\item LLM 使用=True；回退=无；后端=llm\_slot\_frame+runtime\_semantic。
\item 第 3 轮：动作=finalize，阶段=source，闭环状态=True。
\item 缺失字段：无。
\item LLM 使用=True；回退=无；后端=llm\_semantic\_frame+runtime\_semantic+confirmed\_pending\_overwrite。
\end{itemize}
\subsection{案例 MTM09（environmental\_radiation）}
\textbf{对话：}
\begin{itemize}
\item 用户：我想先搭一个空气体积中的低能 gamma 传播场景，细节稍后补。
\item 系统：目前我们已经完成了几何、材料和源的基本设置。接下来还需要补充物理列表、输出格式和输出路径这三项信息。
\item 用户：设为 500 mm x 500 mm x 500 mm 的空气盒，材料 G4\_AIR；点源 gamma，能量 0.5 MeV，位置 (0,0,-300) mm，方向 +z；物理列表 QBBC；输出 root。
\item 系统：我注意到这可能会覆盖已确认的设置。具体修改如下：
- 盒体 x 方向尺寸：1000.0 → 500.0
- 盒体 y 方向尺寸：1000.0 → 500.0

请确认是否应用这些修改：
• 如果确认修改，请直接回复“确认”。
• 如果想保留原值，请直接回复“保持原值”。
\item 用户：确认
\item 系统：配置已完成。当前已确认的设置包括：源位置、源能量、盒体 z 方向尺寸、盒体 y 方向尺寸和盒体 x 方向尺寸。
\end{itemize}
\textbf{内部链路（自然语言）：}
\begin{itemize}
\item 第 1 轮：动作=summarize\_progress，阶段=materials，闭环状态=False。
\item 缺失字段：physics.physics\_list、output.format、output.path。
\item LLM 使用=True；回退=无；后端=llm\_slot\_frame+runtime\_semantic。
\item 第 2 轮：动作=confirm\_overwrite，阶段=source，闭环状态=False。
\item 缺失字段：无。
\item LLM 使用=True；回退=无；后端=llm\_slot\_frame+runtime\_semantic。
\item 第 3 轮：动作=finalize，阶段=physics，闭环状态=True。
\item 缺失字段：无。
\item LLM 使用=True；回退=无；后端=llm\_semantic\_frame+runtime\_semantic+confirmed\_pending\_overwrite。
\end{itemize}
\subsection{案例 MTM10（high\_energy\_physics）}
\textbf{对话：}
\begin{itemize}
\item 用户：Set geometry as a sphere target; I have not decided radius or source details yet.
\item 系统：Here's a summary of your configuration progress:

**Geometry:** You've set up a box-shaped world volume and chosen a sphere as your target. The box dimensions (X and Y) are set, but the sphere's radius is still needed.

**Materials:** You need to assign a material to your sphere volume.

**Source:** The particle source details are still required: its type, the particle to be emitted, and the particle's energy.

**Physics \& Output:** The physics list and the output format and location also need to be specified.
\item 用户：Sphere target radius 80 mm, lead material, electron point source 5 MeV at (0,0,-200) mm toward +z, physics FTFP\_BERT, output hdf5.
\item 系统：Your configuration is complete. The following items have been confirmed: your physics selection reasons, physics selection source, material selection reasons, and material selection source.
\end{itemize}
\textbf{内部链路（自然语言）：}
\begin{itemize}
\item 第 1 轮：动作=summarize\_progress，阶段=geometry，闭环状态=False。
\item 缺失字段：geometry.params.child\_rmax、materials.volume\_material\_map、source.type、source.particle、source.energy、source.position、source.direction、physics.physics\_list、output.format、output.path。
\item LLM 使用=True；回退=无；后端=llm\_slot\_frame+runtime\_semantic。
\item 第 2 轮：动作=finalize，阶段=materials，闭环状态=True。
\item 缺失字段：无。
\item LLM 使用=True；回退=无；后端=llm\_slot\_frame+runtime\_semantic。
\end{itemize}
\subsection{案例 MTM11（materials\_update）}
\textbf{对话：}
\begin{itemize}
\item 用户：Create a complete setup: 100 mm x 100 mm x 100 mm copper box, point gamma 1 MeV at (0,0,-200) mm toward +z, physics QBBC, output json.
\item 系统：Your configuration is complete. I have confirmed the physics and material selections as requested.
\item 用户：Change the target material to G4\_Al and change the output format to root.
\item 系统：This would overwrite an already confirmed value.

Pending changes:
- **Target material**: from Copper (G4\_Cu) to Aluminum (G4\_Al)
- **Output format**: from JSON to ROOT

Reply **"confirm"** to apply these changes, or **"keep existing"** to keep the current values.
\item 用户：confirm
\item 系统：Configuration complete. The following items are now confirmed: volume-to-material binding, output path, output format, and the selected materials.
\end{itemize}
\textbf{内部链路（自然语言）：}
\begin{itemize}
\item 第 1 轮：动作=finalize，阶段=materials，闭环状态=True。
\item 缺失字段：无。
\item LLM 使用=True；回退=无；后端=llm\_slot\_frame+runtime\_semantic。
\item 第 2 轮：动作=confirm\_overwrite，阶段=source，闭环状态=False。
\item 缺失字段：无。
\item LLM 使用=True；回退=无；后端=llm\_slot\_frame+runtime\_semantic。
\item 第 3 轮：动作=finalize，阶段=physics，闭环状态=True。
\item 缺失字段：无。
\item LLM 使用=True；回退=无；后端=llm\_semantic\_frame+runtime\_semantic+confirmed\_pending\_overwrite。
\end{itemize}
\subsection{案例 MTM12（overwrite\_reject）}
\textbf{对话：}
\begin{itemize}
\item 用户：请创建一个完整配置：100 mm x 100 mm x 100 mm 的铜盒，材料 G4\_Cu；点源 gamma，能量 1 MeV，位置 (0,0,-200) mm，方向 +z；物理列表 QBBC；输出 json。
\item 系统：配置已完成。当前已确认的设置包括：物理列表的选择依据和来源，以及材料的选择依据和来源。
\item 用户：把材料改成 G4\_Al，并把输出格式改成 root。
\item 系统：检测到您要修改已确认的设置，具体变更如下：
- 材料：从 G4\_Cu 改为 G4\_Al
- 输出格式：从 json 改为 root

请确认是否应用这些修改？如果确认，请直接回复“确认”；如果希望保留原值，请回复“保持原值”。
\item 用户：保持原值
\item 系统：好的，已为您保留材料的原有值（G4\_Cu）和输出格式的原有值（json）。
\end{itemize}
\textbf{内部链路（自然语言）：}
\begin{itemize}
\item 第 1 轮：动作=finalize，阶段=materials，闭环状态=True。
\item 缺失字段：无。
\item LLM 使用=True；回退=无；后端=llm\_slot\_frame+runtime\_semantic。
\item 第 2 轮：动作=confirm\_overwrite，阶段=source，闭环状态=False。
\item 缺失字段：无。
\item LLM 使用=True；回退=无；后端=llm\_slot\_frame+runtime\_semantic。
\item 第 3 轮：动作=reject\_overwrite，阶段=physics，闭环状态=True。
\item 缺失字段：无。
\item LLM 使用=True；回退=无；后端=llm\_semantic\_frame+runtime\_semantic+rejected\_pending\_overwrite。
\end{itemize}
\end{document}